% !TeX root = ../../../main.tex
\subsubsection{Ablation Study: How ATE Estimates Change When Risk Scores are Removed}
\label{ch03:sec:04_results_ablation}
% \vspace{-0.25em}

We next examine the impact of removing the county-provided risk scores on the ATE estimates for different survival analysis model and causal inference method pairs. 
\ChThreeTableRefs{ch03:tab:table_ablation_CoxPH,ch03:tab:table_ablation_RandomSurvivalForest,ch03:tab:table_ablation_DeepSurv,ch03:tab:table_ablation_DeepHit,ch03:tab:table_ablation_RandomSurvivalForest} in Appendix~\ref{ch03:apd:ablation} compare the estimates obtained with the full model versus those from the ablated model (without risk scores) at the snapshot time 3, 6, 9, and 12 months. 
Notably, the magnitude of the negative ATE increases in the ablated models across all time horizons, suggesting that the risk scores partially capture confounding factors. 
Their removal amplifies the estimated impact of non-adherence, suggesting that these scores encode important latent risk factors that, when included, help account for confounding effects.

% \vspace{-1.5em}
